\documentclass[10pt,a4paper]{amsart}
\usepackage[margin=19mm]{geometry}
\usepackage{amsmath,amssymb,mathtools}
\usepackage[scr=rsfs,cal=euler]{mathalfa}
\usepackage{xfrac}
\usepackage[unicode,hidelinks,bookmarks=false]{hyperref}
\newcommand{\ie}{i.e.\ }
\newcommand{\cf}{cf.\ }
\newcommand{\resp}{respectively\ }
\newcommand{\Subsection}{\subsection}
\providecommand{\nobreakdash}{\nobreak\hbox{-}}
\setlength{\emergencystretch}{2em}
\multlinegap0pt
\usepackage[all]{xy}
\usepackage{enumerate}
\usepackage{bm}
\usepackage{amssymb}
\usepackage{tikz}
\usepackage{mathtools}
\newtheorem{lem}{Lemma}
\DeclareMathOperator{\Li}{Li}
\def\N{\mathbb{N}}
\title{\bf Contour Integrations and Parity Results of Hurwitz-type Cyclotomic Euler Sums}
\author{Hongyuan Rui}
\date{Source: arXiv:2601.00035v1 (CC BY 4.0)}
\begin{document}
\maketitle

\noindent\emph{Source and licence.} \bf Contour Integrations and Parity Results of Hurwitz-type Cyclotomic Euler Sums. Version arXiv:2601.00035v1 (CC BY 4.0), distributed by arXiv under the Creative Commons Attribution 4.0 International licence, http://creativecommons.org/licenses/by/4.0/, as recorded in the arXiv licence metadata for this exact version and shown on the abstract page https://arxiv.org/abs/2601.00035v1. Full licence notice: \texttt{licenses/arxiv-2601.00035-CC-BY-4.0.txt}.\par\smallskip

\noindent\textit{Editorial scope.} Counted unit: the article's lem lem-extend-xuzhou-one (source0.tex lines 385--395). The article's complete original proof of it is delivered with it (source0.tex lines 396--406, one \texttt{proof} environment of 11 lines). No mathematics is written, repaired or re-derived: the statement and the proof are the article's own lines, in the article's own notation, and no retained line is substituted or re-flowed.  No further source-local context is needed: every cross-reference of every retained line resolves inside the two delivered spans.  Nothing was omitted from any retained span, so the package carries no omission record.  No retained line contains a citation, so the wrapper carries no bibliography and no bibliography entry.  No retained line contains a figure environment, an image-inclusion command or drawing code, so no image is delivered and the package declares no asset dependency.  Editorial formatting only, in addition: an AMS \texttt{amsart} preamble that restates the article's own declarations and macros used by the retained lines, namely the environments \texttt{lem} on separate counters, together with the standard packages those lines need.  The retained environments print in this document as Lemma 1 in order of appearance, whereas the article prints the same items with the numbering of its own full document.  The complete native source of the article as distributed in the arXiv e-print for this exact version is bundled unchanged as \texttt{sources/191966/source0.tex}.
\begin{lem}\label{lem-extend-xuzhou-one} For $p\in \N$, if $|s+n|<1\ (n\geq 0)$, then
\begin{align}\label{Lexp-phi--n-diffp-1}
&\frac{\phi^{(p-1)}(s+a;x)}{(p-1)!} (-1)^{p-1}\nonumber\\
&=x^n\sum_{k=0}^\infty \binom{k+p-1}{p-1} \left((-1)^k \Li_{k+p}(x;a)x^{-1}+(-1)^p \zeta_n\Big(k+p;x^{-1};-a\Big)\right)(s+n)^k
\end{align}
and if $|s-n|<1\ (n\geq 1)$
\begin{align}\label{Lexp-phi-n-diffp-1}
&\frac{\phi^{(p-1)}(s+a;x)}{(p-1)!} (-1)^{p-1}\nonumber\\
&=x^{-n-1}\sum_{k=0}^\infty \binom{k+p-1}{p-1} (-1)^k  \left( \Li_{k+p}(x;a)-\zeta_{n}\Big(k+p;x;a-1\Big)\right)(s-n)^k.
\end{align}
\end{lem}

\begin{proof}
If $|s+n|<1\ (n\geq 0)$, it follows directly from the definition that
\begin{align*}
\phi(s+a;x)=x^n \sum_{m=0}^\infty \left((-1)^m \Li_{m+1}(x;a)x^{-1}-\zeta_n\Big(m+1;x^{-1};-a\Big) \right)(s+n)^m.
\end{align*}
Taking the $(p-1)$th derivative with respect to $s$ on both sides of the above equation yields formula \eqref{Lexp-phi--n-diffp-1}. Similarly, if $|s-n|<1\ (n\geq 1)$, by a direct calculation, we obtain
\begin{align*}
\phi(s+a;x)=x^{-n-1} \sum_{m=0}^\infty (-1)^m \left( \Li_{m+1}(x;a)-\zeta_n\Big(m+1;x;a-1\Big) \right)(s-n)^m.
\end{align*}
Taking the $(p-1)$th derivative with respect to $s$ on both sides of the above equation yields formula \eqref{Lexp-phi-n-diffp-1}.
\end{proof}

\end{document}
